\section{Smooth continuous envelopes from the complete BFJ model estimate}
\label{sec:smooth}
\textit{Support: G331 and earlier own n1 (Section~\ref{sec:canonical}). The BFJ envelope estimate is a new published application, not an envelope grant imported from G331.}

In this section only, $S$ is the local ring at a smooth closed point of an affine variety over an algebraically closed characteristic-zero field. We first use its maximal-ideal gauge. Dimension at least two is the exact BFJ setting; the DVR and empty cases were treated directly in n1. Completion and primary contraction are those of G331 and transport functions on every semivaluation, not merely fraction-field valuations.

We apply Boucksom--Favre--Jonsson, \emph{Valuations and plurisubharmonic singularities}, Proposition~3.18, including its entire immediate proof and Lemma~3.22, together with its stated function version Remark~3.21. Definitions~3.15--3.17 and the immediate preceding multiplier-ideal units were read in full. The setting, retractions, closed-simplex continuity and all-retraction extension are precisely the published ones. Smooth principalization and usual multiplier-ideal subadditivity are named deeper foundations; they are not extended to singular or positive-characteristic germs. The reading ledger gives the complete windows and the boundary of this source audit.

Let $F$ be a finite continuous function determined on a single BFJ dual complex: $F=F\circ r_{\pi_0}$. A signed real linear combination of primary ideal values has this property on a common principalization (BFJ~1.14), and these functions are uniformly dense (1.15, also G331's independently proved lattice density). Set $M=\max(0,-\min F)$. The constant $-M$ is a nef primary real profile below $F$.

Consider the weak nef Weil envelope
\[
 W=\sup\{Z:Z\text{ nef},\ Z\le Z(F)\}.
\]
This is a finite nef Weil divisor: restrict the competitors by replacing them with their maximum with the barrier $-M$; maxima preserve nefness and the obstacle bound. They form a directed family with finite coefficient bounds $-M\le Z\le0$. Their componentwise limit lies in the weakly closed nef cone. Thus $-M\le W\le0$ and $W\le Z(F)$. This is the published Definition~3.17 construction, with its bounds made explicit.

By BFJ~3.11 let $H$ be its formal psh function, extended to every semivaluation by
\[
 H(v)=\inf_{\pi}H(r_\pi v).
\]
The inequality $H\le F$ holds on all supports. On every closed simplex it follows from the coefficient inequality and density of its divisorial rational points, using the actual closed-simplex continuity theorem. For $\pi\ge\pi_0$, $r_{\pi_0}r_\pi=r_{\pi_0}$, hence $F(r_\pi v)=F(v)$. Taking the infimum over that cofinal family gives $H(v)\le F(v)$. Also $H\ge-M$ everywhere. Every nonbottom $u\in D$ below $F$ is, by the full qualified smooth comparison in G331, a BFJ formal psh function with nef coefficients below $Z(F)$. Those coefficients are below $W$, and their all-retraction extensions are below $H$. The bottom function is harmless. Thus $H=p_S(F)$ as a raw function, and $H$ is already an actual $D$ envelope. Its continuity is still to be proved.

Put $B_m=L^\infty(mW)$ in the completed ring, using BFJ~3.15. By 3.17,
\[
 L^\infty(mW)=L^\infty(mZ(F)). \tag{7}
\]
These ideals contain the primary power $\widehat{\mathfrak m}^{\lceil mM\rceil}$, with a positive power chosen when necessary. In the nonzero-envelope case they are proper. After exact primary contraction they are exactly the canonical ideals of n1 for $F$.

Here is a complete check of the latter assertion on all supports. The canonical inequality implies the coefficient tests at every divisorial point. Conversely let a formal $g$ pass those coefficient tests, and choose $N\ge mM$, $N>0$. The primary-or-unit ideal $(g)+\widehat{\mathfrak m}^N$ has divisorial profile
\[
 \log|(g)+\widehat{\mathfrak m}^N|=\max(\log|g|,-N)\le mF.
\]
Its profile is a BFJ all-retraction extension and a continuous geometric ideal profile. As in the preceding paragraph, the divisorial inequality extends through the closed simplices and all refinements of $\pi_0$ to every support. Since $v(g)\ge v((g)+\widehat{\mathfrak m}^N)$, it gives $v(g)+mF(v)\ge0$ everywhere, including infinite $v(g)$. The completion homeomorphism of G331 then contracts $B_m$ precisely to $A_m^S(F)$. This argument does not map the original fraction field through a nonminimal support and does not divide an infinite value.

If $W=0$, then $0\le F$ on the dense divisorial points and hence on the whole space, and the envelope and all profiles are zero. Otherwise BFJ~3.14 gives strict negativity of $H$, uniformly bounded above by a negative constant. On the fixed compact dual complex, thinness $A$ is bounded and $H$ is continuous and strictly negative. Choose a finite constant $C>0$ as in Proposition~3.18; its proof allows $C$ with $A\le-CW$ on this single complex. It uses thinness only there, not on the entire boundary.

Write $W_t=t^{-1}Z(L^\infty(tW))$ for positive real $t$. The published inequality 3.18(iii) says
\[
 W_t\le W\le (1-C/t)W_{t-C}\qquad(t>C).
\] At $t=m+C$ this implies
\[
 (1+C/m)W\le W_m\le W.
\]
Every displayed coefficient inequality extends to all semivaluations by closed-simplex continuity and the infimum over all retractions. The extension of $W_m$ is exactly $h_m=-v(A_m^S(F))/m$. Since $-M\le H\le0$, we obtain the uniform, all-support estimate
\[
 0\le H-h_m\le CM/m\qquad(m\ge1). \tag{8}
\]
Thus $H$ is a uniform limit of actual positive primary Cartier profiles and belongs to the continuous independently geometric class $C$. No finite generation of the canonical algebra was used. This is an obstacle-dependent primary approximation obtained from a smooth-model estimate, with every infinity-support endpoint included. It does not use a singular Jacobian correction.

Now let $F$ be any finite continuous signed obstacle on the smooth boundary. Choose signed primary model functions $F_j\to F$ uniformly by the complete density statement. Their envelopes $H_j\in C$ exist by (8). The raw Lipschitz estimate (4) makes $H_j$ uniformly Cauchy and identifies its limit with $p_S(F)$. Uniform closedness of the independently geometric class proves that $p_S(F)\in C$ and is the greatest $D$ minorant. This establishes arbitrary continuous-obstacle regularity at the precise published smooth hypotheses, rather than assuming it as an extra G331 grant.

All primary gauges follow the exact homogeneous transport of n1. If $\nu$ has maximal-ideal normalization and $b$ is another primary gauge, then $w=\nu/\nu(b)$ and $F_{\mathfrak m}(\nu)=\nu(b)F_b(w)$. For a signed primary model obstacle, $F_{\mathfrak m}$ is still a signed real primary model function. Estimate (8) in the $b$ gauge divides by $\nu(b)$, whose positive minimum is finite on the compact section. For arbitrary continuous obstacles the same transport preserves continuity and the greatest-minorant property. None of these statements applies BFJ to a smooth formal germ outside its algebraic closed-point domain: completion transport supplies precisely the already granted function identification, not an enlarged source theorem.
